\part{Solutions}

\section{Part I}
\paragraph{Solution to Activity 1}
\begin{enumerate}
	\item \begin{enumerate}
	\item Recall that the \(i\)-th column of \(AB\), the product of two matrices, is the linear combination of the columns of \(A\) with the entries of the \(i\)-th 
		column vector of \(B\), \(\vec{b}_i\), acting as weights. Another way to put it is that
		\(AB = \begin{bmatrix}A\vec{b}_1 & A\vec{b}_2 & \dots & A\vec{b}_p\end{bmatrix}\), for
		\(A \in \mathbb{R}^{m \times n}\) and \(B \in \mathbb{R}^{n \times p}\). Knowing this, the second column of \(AB\) is \(A\vec{b}_2\).
		\[\begin{bmatrix}1 & 0 & 3 \\ 1 & 5 & 6 \\ 2 & 1 & 1\end{bmatrix}\begin{bmatrix}0 \\ 1 \\ 3\end{bmatrix} 
		  = 0\begin{bmatrix}1 \\ 0 \\ 3\end{bmatrix} + 1\begin{bmatrix}0 \\ 5 \\ 1\end{bmatrix} + 3\begin{bmatrix}3 \\ 6 \\ 1\end{bmatrix}
		  = \begin{bmatrix}9 \\ 23 \\ 4\end{bmatrix}\]
	\item We could use the fact that left-multiplying a matrix by another (appropriate) matrix (say \(AB\)) is similar to taking a linear combination of the
		\emph{rows} of the right matrix (\(B\) in this case) with the entries of the \(i\)-th \emph{row} of \(A\) acting as weights. In this case, the second row of	
		\(AB\) is \(\row_2(A) \cdot B\), or 
		\[\begin{bmatrix}1 & 5 & 6\end{bmatrix}\begin{bmatrix}2 & 0 & -1 \\ 3 & 1 & 1 \\ 0 & 3 & 0\end{bmatrix} 
		   = 1\begin{bmatrix}2 & 0 & -1\end{bmatrix} + 5\begin{bmatrix}3 & 1 & 1\end{bmatrix} + 6\begin{bmatrix}0 & 3 & 0\end{bmatrix}
		   = \begin{bmatrix}17 & 23 & 4\end{bmatrix}\]
	\end{enumerate} 	
	\item Use the fact that columns of \(AB\) are the linear combinations of the columns of \(A\). If 
		\(B = \begin{bmatrix}\vec{b}_1 & \vec{b}_2 & \vec{b}_1 + \vec{b}_2\end{bmatrix}\). Then, 
		\[AB = \begin{bmatrix}A\vec{b}_1 & A\vec{b}_2 & A(\vec{b}_1 + \vec{b}_2)\end{bmatrix} 
		       = \begin{bmatrix}A\vec{b}_1 & A\vec{b}_2 & A\vec{b}_1 + A\vec{b}_2\end{bmatrix}\]
		Notice that the last column of \(AB\) is the sum of the first two columns, meaning that the columns of \(AB\) are linearly \emph{dependent}.
	\item You could use the above reasoning, but a better way to think about this is going back to the homogenous equation \(AB\vec{y} = \vec{0}\). Note that 
		since the columns of \(B\) are linearly \emph{dependent}, the equation \(B\vec{x} = \vec{0}\) has a nontrivial solution. We could set 
		\(\vec{y} = \vec{x}\) to get \(AB\vec{x} = 0\). But notice that \(AB\vec{x} = A(B\vec{x}) = A\vec{0} = \vec{0}.\) The equation \(AB\vec{x} = \vec{0}\)
		has a nontrivial solution, so its columns are linearly \emph{dependent}. So, the answer is \textbf{yes}.
\end{enumerate}

\paragraph{Solution to Activity 2}
\begin{enumerate}
	\item The first few steps are
		\begin{equation*} 
			\begin{split}
				\small
				&\hphantom{{}\overset{\tiny E_{3}}{\sim}{}}
				\left(\!\!\begin{array}{ccc|ccc}1 & 0 & 1 & 1 & 0 & 0\\ 2 & 1 & 0 & 0 & 1 & 0\\ 5 & 1 & 1 & 0 & 0 & 1\end{array}\!\!\right) 
				\overset{\tiny E_{1}}{\sim}
				\left(\!\!\begin{array}{ccc|ccc}1 & 0 & 1 & 1 & 0 & 0\\ 0 & 1 & -2 & -2 & 1 & 0\\ 5 & 1 & 1 & 0 & 0 & 1\end{array}\!\!\right)
				\overset{\tiny E_{2}}{\sim}
				\left(\!\!\begin{array}{ccc|ccc}1 & 0 & 1 & 1 & 0 & 0\\ 0 & 1 & -2 & -2 & 1 & 0\\ 0 & 1 & -4 & -5 & 0 & 1\end{array}\!\!\right) \\
				&\overset{\tiny E_{3}}{\sim}
				\left(\!\!\begin{array}{ccc|ccc}1 & 0 & 1 & 1 & 0 & 0\\ 0 & 1 & -2 & -2 & 1 & 0\\ 0 & 0 & -2 & -3 & -1 & 1\end{array}\!\!\right) \sim \dots
			\end{split}
		\end{equation*}
		We've reached the REF of the matrix by just applying 3 row operations, which are 
		\[E_1 = \begin{pmatrix}1 & 0 & 0 \\ -2 & 1 & 0 \\ 0 & 0 & 1\end{pmatrix} \quad E_2 = \begin{pmatrix}1 & 0 & 0 \\ 0 & 1 & 0 \\ -5 & 0 & 1\end{pmatrix}
		\quad E_3 = \begin{pmatrix}1 & 0 & 0 \\ 0 & 1 & 0 \\ 0 & -1 & 1 \end{pmatrix}.\]
		Anyways, moving on
		\begin{equation*}
			\begin{split}
			\small
			\dots &\sim
			\left(\!\!\begin{array}{ccc|ccc}1 & 0 & 1 & 1 & 0 & 0\\ 0 & 1 & 0 & 1 & 2 & -1\\ 0 & 0 & -2 & -3 & -1 & 1\end{array}\!\!\right)
			\sim 
			\left(\!\!\begin{array}{ccc|ccc}1 & 0 & 1 & 1 & 0 & 0\\ 0 & 1 & 0 & 1 & 2 & -1\\ 0 & 0 & -2 & -3 & -1 & 1\end{array}\!\!\right) \\
			&\sim 
			\left(\!\!\begin{array}{ccc|ccc}1 & 0 & 0 & -1/2 & -1/2 & 1/2\\ 0 & 1 & 0 & 1 & 2 & -1\\ 0 & 0 & 1 & 3/2 & 1/2 & -1/2\end{array}\!\!\right) = A^{-1}
			\end{split}
		\end{equation*}
	\item We've already written out the elementary matrices in the previous part. Now, to combine them, you \emph{could} multiply them out. A simpler way is 
		to realize that multiplying (lower triangular) elementary matrices essentially ``fills" in the lower part of the matrix. That is,
		\[\begin{pmatrix}1 & 0 & 0 \\ \color{magenta} -2 & 1 & 0 \\ 0 & 0 & 1\end{pmatrix}
		   \begin{pmatrix}1 & 0 & 0 \\ 0 & 1 & 0 \\ \color{teal} -5 & 0 & 1\end{pmatrix}
		   \begin{pmatrix}1 & 0 & 0 \\ 0 & 1 & 0 \\ 0 & \color{red}-1 & 1 \end{pmatrix} 
		= \begin{pmatrix}1 & 0 & 0 \\ \color{magenta}-2 & 1 & 0 \\ \color{teal} -5 & \color{red}-1 & 1 \end{pmatrix}\]

	\item \[{\color{teal} AB} = {\color{orange} CA} = {\color{magenta} I_n} \implies C{\color{teal} AB} = C \cdot {\color{magenta}I_n}\implies ({\color{orange}CA}){\color{teal}B} =  C \implies B = C.\] 
		You could also cite the Invertible Matrix Theorem (IMT), stating that each square matrix has
		only one unique inverse.
\end{enumerate}

\paragraph{Solution to Activity 3}
\begin{enumerate}
	\item \textbf{Yes} Every invertible matrix \(A \in \mathbb{R}^{m \times n}\) possesses \(n\) pivots, meaning that there is a pivot in every row and column. It follows
		then that the every invertible matrix's RREF is the identity matrix \(I_n\). You could also just cite IMT as the reasoning.
	\item \textbf{True} Again, you can just cite IMT. The main reasoning here is the fact that the mapping \(\vec{x} \mapsto A\vec{x}\) is onto, therefore the standard
		matrix \(A\) must have \(n\) pivots. This means that it's row equivalent to the corresponding idenity matrix \(I_n\), which means it's invertible by IMT.
	\item \textbf{True} Since the transformation is one-to-one, that means that there's \(n\) pivot positions in the matrix. We have \(n\) rows in the matrix (Why?) and
		so the columns of the matrix span \(\mathbb{R}^n\).
	\item \textbf{False} If \(A \in \mathbb{R}^{n \times n}\) and the columns span \(\mathbb{R}^n\), then \(A\) is invertible. By IMT, \(A^T\) is also invertible, and so
		its columns (which are the rows of \(A\)) span \(\mathbb{R}^n\). So the rows must also span \(\mathbb{R}^n\). (You could also say that there are \(n\) pivot
		rows, so the rows span \(\mathbb{R}^n\) as well).
	\item \textbf{False} A square invertible matrix cannot have multiple inverses. Look at \textbf{Activity 2} (Question 3).
\end{enumerate}

\paragraph{Solution to Activity 4}
\begin{enumerate}
	\item There are multiple ways to solve this problem. I go with the ``standard" approach. We need a matrix 
		\(A^{-1} = \begin{bmatrix}X & Y \\ W & Z \end{bmatrix}\) such that \(AA^{-1} = A^{-1}A = I_{2n}.\) We need
		\[\begin{bmatrix}B & I_n \\ 0 & C \end{bmatrix}\begin{bmatrix}X & Y \\ W & Z \end{bmatrix} = \begin{bmatrix}I_n & 0 \\ 0 & I_n\end{bmatrix}.\]
		If we multiply the left-side out, we get
		\[\begin{bmatrix}BX + W & BY + Z \\ CW & CZ \end{bmatrix}= \begin{bmatrix}I_n & 0 \\ 0 & I_n\end{bmatrix}.\]
		So, we need
		\begin{align*}
			BX + W &= I \implies X = B^{-1}\\
			BY + Z &= 0 \implies Y = -B^{-1}C^{-1}\\
			CW &= 0 \implies W = 0\\
			CZ &= I \implies Z = C^{-1}
		\end{align*}
	\item The dimensions of \(H\) are \(q \times p\), and the dimensions of \(A\) are \((m + q) \times (p + l)\).

\paragraph{Solution to Activity 5}
\begin{enumerate}
	\item Every invertible matrix could be written as the product of elementary matrices. Since \(B\) is invertible, that means \(B = E_p\dots E_3E_2E_1\). Substituting 
		that in, we get that \(A = (E_p\dots E_3E_2E_1)C\). Since \(C\) is left multiplied by a bunch of elementary matrices, we know that there is a sequence of row
		operations that when applied to \(C\) produce \(A\). Hence, \(A \sim C\).
	\item The matrix we found was \(\begin{pmatrix}1 & 0 & 0 \\ \color{magenta}-2 & 1 & 0 \\ \color{teal} -5 & \color{red}-1 & 1 \end{pmatrix}\). The way 
		inverses work for the product of (lower trangular) elementary matrix is that you negate the entries under the main diagonal. (It might be worthwhile to 
		convince yoursellf why.) Therefore, in this case, 
		\[\begin{pmatrix}1 & 0 & 0 \\ \color{magenta}-2 & 1 & 0 \\ \color{teal} -5 & \color{red}-1 & 1 \end{pmatrix}^{-1} 
		= \begin{pmatrix}1 & 0 & 0 \\ \color{magenta}2 & 1 & 0 \\ \color{teal} 5 & \color{red}1 & 1 \end{pmatrix}.\] The \(LU\)-factorization of the matrix 
		\(A\) is therefore \[\begin{pmatrix}1 & 0 & 1 \\ 2 & 1 & 0 \\ 5 & 1 & 1\end{pmatrix} = \begin{pmatrix}1 & 0 & 0 \\ \color{magenta}2 & 1 & 0 \\ \color{teal} 5 & \color{red}1 & 1 \end{pmatrix}\begin{pmatrix}1 & 0 & 1 \\ 0 & 1 & -2 \\ 0 & 0 & -2 \end{pmatrix}\] 

	\item We need \(\vec{x}\) such that \(A\vec{x} = L\underbrace{(U\vec{x})}_{\vec{y}} = L\vec{y} = \begin{pmatrix}1 \\ 3 \\ 1\end{pmatrix}.\) We can now solve
		for the simpler system \(L\vec{y} = \begin{pmatrix}1 \\ 3 \\ 1\end{pmatrix}\) first then solve for \(U \vec{x} = \vec{y}.\)

		For \(\vec{y}\):
			\[\left(\!\!\begin{array}{ccc|c}1 & 0 & 0 & 1 \\ 2 & 1 & 0 & 3 \\ 5 & 1 & 1 & 1\end{array}\!\!\right) 
			  \sim
			  \left(\!\!\begin{array}{ccc|c}1 & 0 & 0 & 1 \\ 0 & 1 & 0 & 1 \\ 0 & 1 & 1 & -4\end{array}\!\!\right)
			  \sim
			  \left(\!\!\begin{array}{ccc|c}1 & 0 & 0 & 1 \\ 0 & 1 & 0 & 1 \\ 0 & 0 & 1 & -5\end{array}\!\!\right)				
			\]
		So, \(\vec{y} = \begin{pmatrix}1 \\ 1 \\ -5\end{pmatrix}\). Now, to find \(\vec{x}\), solve the system \(U\vec{x} = \vec{y}.\)
			\[\left(\!\!\begin{array}{ccc|c}1 & 0 & 1 & 1 \\ 0 & 1 & -2 & 1 \\ 0 & 0 & -2 & -5 \end{array}\!\!\right) 
			  \sim
			   \left(\!\!\begin{array}{ccc|c}1 & 0 & 1 & 1 \\ 0 & 1 & -2 & 1 \\ 0 & 0 & -2 & -5 \end{array}\!\!\right)
			  \sim \dots \sim
			  \left(\!\!\begin{array}{ccc|c}1 & 0 & 0 & -3/2 \\ 0 & 1 & 0 & 6 \\ 0 & 0 & 1 & 5/2 \end{array}\!\!\right)
			\]
		So, \(\vec{x} = \begin{pmatrix}-3/2 \\ 6 \\ 5/2\end{pmatrix}\).
\end{enumerate}
\end{enumerate}

\section{Part II}

\paragraph{Solution to Activity 6}
\begin{enumerate}
	\item \textbf{No} It's important to realize that by ``line," we mean the subset of vectors of \(\mathbb{R}^n\) such that the ``tip" of the vector touches the line.
				  Anyhow, sometimes we have ``lines" in \(\mathbb{R}^n\) that are not subspaces (such as the solutionset of an inhomogenous system.) The line
				  must pass through the origin in order to be a subspace.
	\item \textbf{Yes} \(\vec{0} \in \vspan\{\vec{v}_1, \dots, \vec{v}_p\}\), and the sum of any two vectors in \(\vspan \{\vec{v}_1, \dots, \vec{v}_p\}\) lies in the 
			          span of the set of vectors. So does the scalar multiple. It's worthwhile to check these facts for yourself.
	\item \textbf{Yes} \[A(2\vec{x} - 2\vec{y}) = 2A\vec{x} - 2A\vec{y} = 2\vec{b} - 2\vec{b} = \vec{0}.\]
	\item \[\begin{bmatrix}1 & 6 & 4 \\ 1 & 3 & 0\end{bmatrix} \sim \begin{bmatrix}1 & 6 & 4 \\ 0 & -3 & -4\end{bmatrix}.\] Since there's a pivot in the first two 
		columns, the first two columns of \(\begin{bmatrix}1 & 6 & 4 \\ 1 & 3 & 0\end{bmatrix}\) form \emph{a} basis for 
		\(\colspace \begin{bmatrix}1 & 6 & 4 \\ 1 & 3 & 0\end{bmatrix}\).
		
		Continuing off, 
		\[\begin{bmatrix}1 & 6 & 4 \\ 0 & -3 & -4\end{bmatrix} 
		  \sim \begin{bmatrix}1 & 0 & -4 \\ 0 & -3 & -4\end{bmatrix}
		  \sim \begin{bmatrix}1 & 0 & -4 \\ 0 & 1 & 4/3\end{bmatrix}\]
		This corresponds to the system
		 \begin{align*}
			      x_1 &        - 4x_3       = 0 \implies x_1 = 4x_3 \\
			      x_2 &+\frac{4}{3}x_3 = 0 \implies x_2 = -\frac{4}{3}x_3
		\end{align*}
		So 
		\[\vec{x} = \begin{pmatrix}x_1 \\ x_2 \\ x_3\end{pmatrix} = \begin{pmatrix}4x_3 \\ -4/3x_3 \\ x_3\end{pmatrix} = x_3\begin{pmatrix}4 \\ -4/3 \\ 1\end{pmatrix}.\]
		Notice what this is saying. It's saying that every \(\vec{x} \in \nullspace\biggl(\begin{bmatrix}1 & 6 & 4 \\ 1 & 3 & 0\end{bmatrix}\biggr)\) is a scalar multiple of
		\(\begin{pmatrix}4 \\ -4/3 \\ 1\end{pmatrix}\). That means that \(\begin{pmatrix}4 \\ -4/3 \\ 1\end{pmatrix}\) spans the null space of the matrix, and since it's linearly
		independent, it is \emph{a} basis for the nullspace of this matrix.

	\item \textbf{No!} For example, in the previous part of this activity, we found that the basis of the column space of \(\begin{bmatrix}1 & 6 & 4 \\ 1 & 3 & 0\end{bmatrix}\)
				   is \(\biggl\{\begin{bmatrix}1 \\ 1\end{bmatrix}, \begin{bmatrix}6\\ 3\end{bmatrix}\biggr\}\). Notice that the matrix had three columns, but the basis has
				   only two. 
\end{enumerate}

\paragraph{Solution to Activity 7}
\begin{enumerate}
	\item Since a basis for the matrix's column space has two vectors in it, that means the column space has a dimension of 2. Since the null space's basis has one vector in it, it's
		of dimension 1. Another way to find the dimension of the null space was to use the Rank-Nullity Theorem that states that the sum of the dimension of these two 
		subspaces equals the number of the columns the matrix has.
	\item For reference, the matrix is \(\begin{pmatrix}1 & 0 & 1 \\ 2 & 1 & 0 \\ 5 & 1 & 1\end{pmatrix}.\) The basis formed is therefore 
		\[\mathcal{B} = \left\{\!\begin{pmatrix}1 \\ 2 \\ 5\end{pmatrix}, \begin{pmatrix}0 \\ 1 \\ 1\end{pmatrix}, \begin{pmatrix}1 \\ 0 \\ 1\end{pmatrix}\!\right\}.\]
		Since \(\mathcal{B}\) forms a basis, we know that we could write the vector \(\begin{pmatrix}1 \\ 2 \\ 1\end{pmatrix}\) as a linear combination of the basis vectors.
\[\begin{pmatrix}1 \\ 2 \\ 1\end{pmatrix} = c_1\begin{pmatrix}1 \\ 2 \\ 5\end{pmatrix} + c_2\begin{pmatrix}0 \\ 1 \\ 1\end{pmatrix} + c_3\begin{pmatrix}1 \\ 0 \\ 1\end{pmatrix} \implies \begin{pmatrix}1 & 0 & 1 \\ 2 & 1 & 0 \\ 5 & 1 & 1\end{pmatrix}\begin{pmatrix}c_1 \\ c_2 \\ c_3\end{pmatrix} = \begin{pmatrix}1 \\ 2 \\1 \end{pmatrix}\]

		Note \(\left[\!\begin{pmatrix}1 \\ 2 \\1 \end{pmatrix}\!\right]_{\mathcal{B}} = \begin{pmatrix}c_1 \\ c_2 \\ c_3\end{pmatrix}\). Solving the inhomogenous system
		is equivalent to finding the \(\mathcal{B}\)-coordinates of the vector.
		\[\left(\!\!\begin{array}{ccc|c}1 & 0 & 1 & 1\\ 2 & 1 & 0 & 2\\ 5 & 1 & 1 & 1\end{array}\!\!\right) \sim \dots \sim \left(\!\!\begin{array}{ccc|c}1 & 0 & 0 & -1\\ 0 & 1 & 0 & 4\\ 0 & 0 & 1 & 2\end{array}\!\!\right).\]
		So, 
		\(\left[\!\begin{pmatrix}1 \\ 2 \\1 \end{pmatrix}\!\right]_{\mathcal{B}} = \begin{pmatrix}c_1 \\ c_2 \\ c_3\end{pmatrix} = \begin{pmatrix}-1 \\ 4 \\ 2\end{pmatrix}.\)
\end{enumerate}